\chapter{Certify, don't identify: a theory of safety without identification}\label{ch:p3}

This chapter is the third paper of the thesis and the candidate's own work. It places the results of Chapters~\ref{ch:p1} and~\ref{ch:p2} in one framework, proves when a question about a market's unknown agents is cheap or impossible to certify, and extends the safety results to agents that are nonlinear and keep learning. Results proved here are marked \Proved{}; results proved in Chapters~\ref{ch:p1}--\ref{ch:p2} are marked accordingly; the conjecture is marked as such.

\begin{quote}\small\textbf{Summary.} Regulators ask to know which algorithms trade in a market; we ask which questions about those algorithms can be answered from market data at all, and at what cost. Every question an outside observer can ask is a functional of the agents' merged measure of memory rates and weights, observed through noisy moments. For linear questions we determine the number of experiments needed to certify the answer, up to a constant factor in accuracy, by a modulus of continuity whose dual is a problem of polynomial approximation with penalized coefficients. A question is certifiable at the parametric rate exactly when its kernel has square-summable power-series coefficients; questions singular at infinite memory cannot be certified without a memory bound. Safety depends on the agents only through their total feedback weight, which sits on the boundary: it costs exactly $\sigma^2\epsilon^{-2}\log(1/\epsilon)$ experiments, and the logarithm cannot be removed, whereas identifying the agents costs polynomially or logarithmically more. The summed-gain guard remains safe for agents whose rules are nonlinear and whose memories and weights change over time, given only a bound on their total sensitivity, and when agents keep learning a certificate lasts for its safety margin divided by their learning speed. The safety bound extends to several assets with positive cross-impact, sharply. A discontinuous guard can fail to admit a strategic equilibrium.\end{quote}

\section{Markets with adaptive agents}\label{sec:fw_model}

\begin{definition}[Cohort economy]\label{def:cohort}
A state $x_t\in\R$, the displacement of the price from its fundamental value, evolves as
\begin{equation}\label{eq:fw_plant}
x_{t+1}=a\,x_t+\kappa\,(u_t+w_{t+1}),\qquad u_t=\sum_{i=1}^K q_i\,\varphi_i(y_{i,t}),\qquad y_{i,t}=r_i\,y_{i,t-1}+(1-r_i)\,x_t,
\end{equation}
with persistence $a\in[0,1)$, impact $\kappa>0$, exogenous shocks $w_t$, and $K$ cohorts of agents with weights $q_i\ge0$, memory rates $r_i\in[0,1)$ and response rules $\varphi_i$. The economy is \emph{linear} when $\varphi_i(y)=y$ for all $i$.
\end{definition}

The exponential state $y_{i,t}$ is the reduced form of steady-state Kalman filtering and of regularized online gradient descent (Lemma~\ref{p1:lem:micro}); Section~\ref{sec:bl_robust} relaxes linearity and stationarity. In the linear economy the aggregate response is a linear filter of past states,
\[
u_t=\sum_{m\ge0}k_m\,x_{t-m},\qquad k_m=\int r^m\,dM(r),\qquad M=\sum_{i=1}^K q_i(1-r_i)\,\delta_{r_i},
\]
so the agents enter the observable dynamics only through the \emph{merged measure} $M$, a positive measure on $[0,1)$, and through its moments $k_m$. Cohorts that share a memory rate are indistinguishable, whatever their identities (Theorem~\ref{p1:thm:merged}). We write $\mathcal M(\bar r,B)$ for the positive measures on $[0,\bar r]$ of total mass at most $B$; $\bar r<1$ is a \emph{memory bound}.

\section{Questions are functionals of the merged measure}\label{sec:fw_questions}

Every question an outside observer can answer from the observable dynamics is a functional $F$ of $M$ (and of the known plant parameters $a,\kappa$). Three families matter.

\paragraph{Identification.} $F(M)=M$ itself, in a metric such as the Wasserstein distance $W_1$, or the atoms $(r_i,q_i)$ when $M$ has $K$ atoms. Chapter~\ref{ch:p1} studies these.

\paragraph{Safety.} For bounded shocks $|w_t|\le\bar w$, the worst-case transient loss of a position $I$ is $|I|\,\bar w\,\|h\|_1$, where $h$ is the closed-loop impulse response, and a pre-trade guard on $\|h\|_1$ admits every safe action and no unsafe one (Lemma~\ref{p2:thm:guard}). In the linear economy with $a\ge0$ all $h_t$ are non-negative, so $\|h\|_1$ equals the summed response
\begin{equation}\label{eq:fw_S}
S(M)=\frac{\kappa}{1-a-\kappa\,\Phi(M)},\qquad \Phi(M)=\sum_{m\ge0}k_m=\int\frac{dM(r)}{1-r}=\sum_i q_i ,
\end{equation}
whenever $a+\kappa\Phi<1$ (Theorem~\ref{p2:thm:sum}). Safety therefore depends on the agents only through one linear functional of $M$, the total feedback weight $\Phi$.

\paragraph{Other operational questions.} The response at a given lag, $k_m=\int r^m dM$; the discounted feedback $\int (1-\beta r)^{-1}dM$; the peak of the impulse response. The first two are linear in $M$; the peak is not, and a guard built on it admits unsafe actions (Chapter~\ref{ch:p2}).

\section{Observation contracts and the cost of a question}\label{sec:fw_contracts}

\begin{assumption}[Reset contract]\label{ass:reset}
The observer can run $N$ independent reset experiments. Each returns the first $L$ aggregate response taps with independent Gaussian noise, $\hat k_m=k_m+\sigma\xi_m$, $\xi_m\sim N(0,1)$, $m=0,\dots,L-1$.
\end{assumption}

This is the controlled-probe contract with observed aggregate flow (contract $I_1$ of Chapter~\ref{ch:p1}). Chapter~\ref{ch:p2} works with the closed-loop taps $h_t$ instead; for the summed response the two contracts give the same rate (Theorem~\ref{p2:thm:certcost}). Correlated tap noise, passive data and tagged disclosure are other contracts, treated in Chapters~\ref{ch:p1} and~\ref{ch:p2}.

\begin{definition}[Cost of certification]\label{def:ncert}
For a functional $F$, a class $\mathcal M$, accuracy $\epsilon>0$ and level $1-\alpha$, $N_{\rm cert}(\epsilon;F,\mathcal M)$ is the smallest $N$ for which some interval built from the data has half-width at most $\epsilon$ and covers $F(M)$ with probability at least $1-\alpha$ for every $M\in\mathcal M$.
\end{definition}

\section{The certification--identification dichotomy}\label{sec:fw_dich}

Consider a linear question $F_g(M)=\int g\,dM$, where $g$ is continuous on $[0,\bar r]$; in the examples $g(r)=\sum_{t\ge0}a_t r^t$ converges absolutely there. Write $k_L(M)=(k_0,\dots,k_{L-1})$ for the first $L$ taps and $p_c(r)=\sum_{t<L}c_tr^t$ for the polynomial with coefficient vector $c\in\R^L$. Truncating the Taylor series gives a first measure of how hard the question is. Its \emph{tail at depth} $T$ and its \emph{coefficient energy} are
\[
\tau_g(T)=B\sup_{0\le r\le\bar r}\Bigl|\sum_{t\ge T}a_t r^t\Bigr|,\qquad A_g(T)=\sum_{t<T}a_t^2,
\]
and its \emph{truncation depth} is $T_g(\epsilon)=\min\{T\ge1:\tau_g(T)\le\epsilon/2\}$; $E_T(g)=\min_{\deg p<T}\max_{[0,\bar r]}|g-p|$ is the error of the best polynomial approximation of degree below $T$. The quantity that decides the cost exactly is the \emph{modulus of continuity} of the question with respect to the observed taps,
\[
\omega_L(\delta)=\sup\bigl\{F_g(M_1)-F_g(M_0):\ M_0,M_1\in\mathcal M(\bar r,B),\ \|k_L(M_1)-k_L(M_0)\|_2\le\delta\bigr\},
\]
the largest difference in the answer between two economies whose observable taps differ by at most $\delta$. For a function $h$ on $[0,\bar r]$ put $\psi(h)=(\sup h)_++(-\inf h)_+$; $\psi(h)\ge\|h\|_\infty$, with equality when $h$ changes sign.

\begin{theorem}[Certification of linear questions]\label{thm:dich}
Under Assumption~\ref{ass:reset}, with $z_\beta$ the standard normal $\beta$-quantile and $\alpha\in(0,1/2)$:
\begin{enumerate}[label=(\roman*)]
\item \emph{Upper bound.} If $L\ge T_g(\epsilon)$, then $N_{\rm cert}(\epsilon;F_g,\mathcal M(\bar r,B))\le\bigl\lceil4z_{1-\alpha/2}^2\sigma^2A_g(T_g(\epsilon))\,\epsilon^{-2}\bigr\rceil$. \hfill\Proved
\item \emph{Truncation is necessary.} For every $T$ there are $M,M'\in\mathcal M(\bar r,B)$ whose first $T$ taps coincide and with $|F_g(M)-F_g(M')|\ge2B\,E_T(g)$. Hence any certificate that uses only the first $T$ taps has worst-case half-width at least $B\,E_T(g)$. \hfill\Proved
\item \emph{Long memory makes certification impossible.} Without a memory bound, that is on the positive measures on $[0,1)$ of mass at most $B$, if $\limsup_{r\uparrow1}|g(r)|\sqrt{1-r}=\infty$ then for every level $1-\alpha$ with $\alpha<1/4$ no finite $N$ certifies $F_g$ to any finite accuracy. \hfill\Proved
\item \emph{Exact cost.} For every depth $L$, the modulus has the dual form
\[
\omega_L(\delta)=\min_{c\in\R^L}\bigl\{\delta\|c\|_2+B\,\psi(g-p_c)\bigr\}\qquad(\delta>0),
\]
and $\omega_L(0)=2B\,E_L(g)$. If $N$ experiments certify $F_g$ to half-width $\epsilon$ at level $1-\alpha$, then $\omega_L(2z_{1-\alpha}\sigma/\sqrt N)\le2\epsilon$. Conversely, if $\omega_L(2z_{1-\alpha/2}\sigma/\sqrt N)\le2\epsilon$, an affine interval $\sum_tc_t\bar k_t+b\pm\epsilon$ certifies $F_g$ with $N$ experiments. Writing $N_\omega(\epsilon)=\min\{N:\omega_L(2z_{1-\alpha}\sigma/\sqrt N)\le2\epsilon\}$ and $\lambda_\alpha=z_{1-\alpha/2}/z_{1-\alpha}$,
\[
N_\omega(\epsilon)\;\le\;N_{\rm cert}(\epsilon;F_g,\mathcal M(\bar r,B))\;\le\;N_\omega(\epsilon/\lambda_\alpha).
\]
\hfill\Proved
\end{enumerate}
\end{theorem}

Part (iv) replaces the matching lower bound in terms of coefficient energy that earlier drafts stated as a target. That statement is false (Remark~\ref{rem:energy}); the modulus is the right quantity, and it determines the cost of every linear question up to the factor $\lambda_\alpha$ in accuracy ($\lambda_{0.05}\approx1.19$). The structure of part (iv), a modulus that governs both a lower bound and an affine interval, is the classical theory of estimating a linear functional over a convex class \citep{t:ibragimov1984,t:donoho1994}. What is specific to markets is the dual form: the cost of a question about the agents is the value of a problem of polynomial approximation of $g$ in which the polynomial's coefficients are penalized, because each coefficient multiplies a noisy tap.

\begin{proof}[Proof of (i)--(iii)]
(i) Let $\bar k_t$ be the average of the $N$ replicates of tap $t$ and $T=T_g(\epsilon)$. The estimator $\hat F=\sum_{t<T}a_t\bar k_t$ has mean $\sum_{t<T}a_tk_t$ and variance $\sigma^2A_g(T)/N$. Its bias is $|\sum_{t\ge T}a_tk_t|=|\int\sum_{t\ge T}a_tr^t\,dM(r)|\le\tau_g(T)\le\epsilon/2$ for every $M\in\mathcal M(\bar r,B)$. The interval $\hat F\pm(\epsilon/2+z_{1-\alpha/2}\sigma\sqrt{A_g(T)/N})$ therefore covers $F_g(M)$ with probability at least $1-\alpha$, and its half-width is at most $\epsilon$ once $z_{1-\alpha/2}\sigma\sqrt{A_g(T)/N}\le\epsilon/2$.

(ii) By duality for best uniform approximation, $E_T(g)=\max\{\int g\,d\nu\}$ over signed measures $\nu$ on $[0,\bar r]$ with total variation at most one that annihilate every polynomial of degree below $T$. Take a maximizer and write $\nu=\nu^+-\nu^-$. Orthogonality to constants gives $\nu^+([0,\bar r])=\nu^-([0,\bar r])=:c\le1/2$. Set $M=(B/c)\nu^+$ and $M'=(B/c)\nu^-$: both lie in $\mathcal M(\bar r,B)$, their first $T$ moments coincide, and $F_g(M)-F_g(M')=(B/c)E_T(g)\ge2B\,E_T(g)$. Two measures that produce identical data cannot both be covered by an interval of half-width below half their distance.

(iii) Compare $M_0=0$ with $M_1=c\,\delta_r$. Under $M_1$ the taps are $k_t=c\,r^t$, so the Kullback--Leibler divergence between the two data distributions over any number of taps is at most $Nc^2/(2\sigma^2(1-r^2))$. Choose $c=\sigma\sqrt{(1-r^2)/N}$, which is at most $B$ for $r$ near one; the divergence is then at most $1/2$, so by Pinsker's inequality the total variation distance between the two data distributions is at most $1/2$. The functional differs by $c\,|g(r)|=\sigma\sqrt{(1-r^2)/N}\,|g(r)|\ge\sigma\sqrt{(1-r)/N}\,|g(r)|$, which is unbounded along a sequence $r\uparrow1$ by assumption. Suppose an interval of half-width $\epsilon$ covered $F_g$ with probability at least $1-\alpha$ under both economies. Once the two values differ by more than $2\epsilon$, the event that the interval contains $F_g(M_0)$ has probability at least $1-\alpha$ under $M_0$ and at most $\alpha$ under $M_1$, so the total variation distance is at least $1-2\alpha>1/2$, a contradiction.
\end{proof}

\begin{proof}[Proof of (iv)]
\emph{Duality.} For any $c\in\R^L$ and any admissible pair, with $h=g-p_c$,
\[
F_g(M_1)-F_g(M_0)=\int h\,dM_1-\int h\,dM_0+\langle c,k_L(M_1)-k_L(M_0)\rangle\le B\psi(h)+\delta\|c\|_2,
\]
because $\int h\,dM_1\le B(\sup h)_+$ and $-\int h\,dM_0\le B(-\inf h)_+$ for measures of mass at most $B$. So $\omega_L(\delta)$ is at most the minimum. For the reverse inequality write $\mu=(M_0,M_1)$ and
\[
\Lambda(\mu,c)=\int h\,d(M_1-M_0)+\delta\|c\|_2,\qquad h=g-p_c .
\]
Since $\inf_c\{\delta\|c\|_2-\langle c,v\rangle\}$ equals $0$ if $\|v\|_2\le\delta$ and $-\infty$ otherwise, $\omega_L(\delta)=\sup_\mu\inf_c\Lambda(\mu,c)$. The set $\mathcal M(\bar r,B)^2$ is convex and compact in the weak-$*$ topology, $\Lambda$ is affine and weak-$*$ continuous in $\mu$ (because $h$ is continuous) and convex and continuous in $c$, so Sion's minimax theorem \citep{t:sion1958} gives $\omega_L(\delta)=\inf_c\sup_\mu\Lambda(\mu,c)$. The inner supremum is $B\psi(g-p_c)+\delta\|c\|_2$, attained by putting mass $B$ at a maximizer of $h$ (if $\sup h>0$) and at a minimizer (if $\inf h<0$). For $\delta>0$ the objective is convex, continuous and coercive in $c$, so the infimum is a minimum. At $\delta=0$ the same argument gives $\omega_L(0)=\inf_c B\psi(g-p_c)$; since adding a constant to $p_c$ can make $g-p_c$ change sign and then $\psi=\sup-\inf$, the infimum is $B$ times the smallest oscillation of $g-p$, that is $2B\,E_L(g)$, which is part (ii) with equality.

\emph{Lower bound.} The averages $\bar k=(\bar k_0,\dots,\bar k_{L-1})$ are sufficient and distributed as $N(k_L(M),(\sigma^2/N)I_L)$, so the total variation distance between the data distributions of two economies is $2\Phi_{\mathcal N}(\sqrt N\|\Delta k\|_2/(2\sigma))-1$, with $\Phi_{\mathcal N}$ the normal distribution function. Take an admissible pair with $\|k_L(M_1)-k_L(M_0)\|_2<2z_{1-\alpha}\sigma/\sqrt N$; the distance is then below $1-2\alpha$. If the two answers differed by more than $2\epsilon$, an interval of half-width $\epsilon$ covering each with probability $1-\alpha$ would separate the two data distributions with total variation at least $1-2\alpha$, as in (iii). Hence $\omega_L(t)\le2\epsilon$ for every $t<2z_{1-\alpha}\sigma/\sqrt N$. The modulus is concave (the admissible pairs form a convex set and the constraint is a norm), finite and nondecreasing, hence continuous on $(0,\infty)$, so $\omega_L(2z_{1-\alpha}\sigma/\sqrt N)\le2\epsilon$.

\emph{Upper bound.} Put $\delta=2z_{1-\alpha/2}\sigma/\sqrt N$, let $c$ attain the minimum in the dual form and $h=g-p_c$. The interval $\sum_tc_t\bar k_t+b\pm\epsilon$ with $b=B[(\sup h)_+-(-\inf h)_+]/2$ has bias at most $B\psi(h)/2$ for every economy, by the bounds in the duality step, and its noise is normal with standard deviation $\sigma\|c\|_2/\sqrt N$. It covers $F_g$ with probability at least $1-\alpha$ whenever $\epsilon\ge B\psi(h)/2+z_{1-\alpha/2}\sigma\|c\|_2/\sqrt N=\omega_L(\delta)/2$.

\emph{Sandwich.} The left inequality restates the lower bound. Concavity and $\omega_L(0)\ge0$ give $\omega_L(\lambda t)\le\lambda\,\omega_L(t)$ for $\lambda\ge1$. If $\omega_L(2z_{1-\alpha}\sigma/\sqrt N)\le2\epsilon/\lambda_\alpha$, then $\omega_L(2z_{1-\alpha/2}\sigma/\sqrt N)\le2\epsilon$, and the upper bound applies.
\end{proof}

\paragraph{What the theorem says.} The cost of answering a linear question is the cost of approximating $g$, on the memory range $[0,\bar r]$, by a polynomial whose coefficients are cheap to estimate. Truncating the Taylor series is one such polynomial and gives part (i); it is often not the best one. Questions singular at infinite memory, faster than $(1-r)^{-1/2}$, cannot be certified without a memory bound (part (iii)). The number of agents and the spacing of their memories do not enter. Two consequences make the picture precise.

\begin{proposition}[Which questions are parametric]\label{prop:parametric}
Let $\omega_\infty(\delta)=\inf_L\omega_L(\delta)$, the modulus when enough taps are observed. Then $\sup_{0<\delta\le1}\omega_\infty(\delta)/\delta<\infty$, so that $N_{\rm cert}=O(\sigma^2\epsilon^{-2})$ for observation depths $L$ large enough, if and only if $g$ agrees on $[0,\bar r]$ with a power series $\sum_ta_tr^t$ whose coefficients are square summable. When $g$ is given by a power series on $[0,\bar r]$ with $\bar r>0$, this is the condition $\sum_ta_t^2<\infty$ on its Taylor coefficients. \hfill\Proved
\end{proposition}

\begin{proof}
If $a\in\ell^2$, take $c=(a_0,\dots,a_{L-1})$: by Cauchy--Schwarz $\psi(g-p_c)\le2\sup_{r\le\bar r}|\sum_{t\ge L}a_tr^t|\le2\|a\|_2\,\bar r^L/\sqrt{1-\bar r^2}$, so $\omega_L(\delta)\le\delta\|a\|_2+2B\|a\|_2\bar r^L/\sqrt{1-\bar r^2}$ and $\omega_\infty(\delta)\le\delta\|a\|_2$. Conversely, suppose $\omega_\infty(\delta)\le C\delta$ for $\delta\in(0,1]$. For $\delta_n=1/n$ choose $L_n$ and $c^{(n)}\in\R^{L_n}$ with $\delta_n\|c^{(n)}\|_2+B\psi(g-p_{c^{(n)}})\le2C\delta_n$. Then $\|c^{(n)}\|_2\le2C$ and $\|g-p_{c^{(n)}}\|_\infty\le2C\delta_n/B\to0$. Padded with zeros, the $c^{(n)}$ are bounded in $\ell^2$, so a subsequence converges weakly to some $a\in\ell^2$. For each $r\in[0,\bar r]$ the sequence $(r^t)_t$ lies in $\ell^2$, hence $p_{c^{(n)}}(r)=\langle c^{(n)},(r^t)\rangle\to\sum_ta_tr^t$ along the subsequence, while $p_{c^{(n)}}(r)\to g(r)$. Uniqueness of power series that agree on an interval gives the last sentence.
\end{proof}

The square-summable power series are the restrictions to $[0,\bar r]$ of the Hardy space $H^2$ of the unit disc, whose functions grow at most like $(1-r)^{-1/2}$ as $r\uparrow1$. That is exactly the growth beyond which part (iii) rules out certification without a memory bound: below the threshold a memory bound makes a question parametric or nearly so, above it nothing helps.

The single tap $g(r)=r^m$ and the discounted feedback $g(r)=(1-\beta r)^{-1}$ with $\beta<1$ have square-summable coefficients and are parametric. The total feedback $\Phi$, the quantity that safety depends on, has $a_t\equiv1$: it sits exactly on the boundary, and it pays exactly one logarithm.

\begin{theorem}[The safety question costs exactly one logarithm]\label{thm:philog}
Let $g(r)=1/(1-r)$, so that $F_g=\Phi$. There are constants $0<c\le C$ and $\delta_0\in(0,1)$, depending only on $(\bar r,B)$, such that for $\delta\le\delta_0$
\[
c\,\delta\sqrt{\log(1/\delta)}\;\le\;\omega_L(\delta)\quad\text{for every }L,\qquad
\omega_L(\delta)\;\le\;C\,\delta\sqrt{\log(1/\delta)}\quad\text{for }L\ge\log(1/\delta)/\log(1/\bar r).
\]
Consequently $N_{\rm cert}(\epsilon;\Phi,\mathcal M(\bar r,B))\asymp\sigma^2\epsilon^{-2}\log(1/\epsilon)$ as $\epsilon\to0$, for deep enough observation, and the same holds for the summed response $S=\kappa/(1-a-\kappa\Phi)$ whenever $a+\kappa B/(1-\bar r)<1$. The logarithmic factor in the certification cost of Chapter~\ref{ch:p2} cannot be removed. \hfill\Proved
\end{theorem}

The lower bound rests on a lemma about polynomials: a polynomial that approximates $1/(1-r)$ closely on the memory range must have large coefficients, because $1/(1-z)$ has a pole at $z=1$ while a polynomial with coefficient norm $R$ grows at most like $R(1-|z|)^{-1/2}$ inside the disc.

\begin{lemma}\label{lem:hardy}
There is $K<\infty$ depending only on $\bar r$ such that for every polynomial $p(r)=\sum_{t<L}c_tr^t$ with $\|c\|_2\le R$,
\[
\log\frac1\eta\le K\max\{R^2,\,(1-\bar r)^{-1}\},\qquad \eta=\sup_{0\le r\le\bar r}\Bigl|\frac1{1-r}-p(r)\Bigr| .
\]
\end{lemma}

\begin{proof}
Put $e(z)=1/(1-z)-p(z)$, $E=[0,\bar r]$, $\rho_1=(1+\bar r)/2$ and $\rho_2=(3+\bar r)/4$. Let $s=\min\{1/(4R^2),(1-\bar r)/4\}$, $x=1-s$ and $\rho=1-s/2$, so that $\rho_1<x<\rho$ and $\rho\ge\rho_2$. By Cauchy--Schwarz $|p(z)|\le R(1-|z|^2)^{-1/2}\le R(1-|z|)^{-1/2}$ in the disc. At $z=x$ this is at most $R/\sqrt s\le1/(2s)$, so $|e(x)|\ge1/s-1/(2s)=1/(2s)$. On the circle $|z|=\rho$, $|e(z)|\le2/s+R\sqrt{2/s}\le3/s$. The function $\log|e|$ is subharmonic on the slit disc $D_\rho\setminus E$, at most $\log\eta$ on $E$ and at most $\log(3/s)$ on $|z|=\rho$, so by the two-constants theorem \citep[Theorem~4.3.7]{t:ransford1995}
\[
\log\frac1{2s}\le\log|e(x)|\le\varpi\log\eta+(1-\varpi)\log\frac3s,
\]
where $\varpi$ is the harmonic measure of $E$ in $D_\rho\setminus E$ seen from $x$. Rearranging, $\varpi\log(1/\eta)\le\log6$. To bound $\varpi$ from below, let $m_*>0$ be the minimum over $|z|=\rho_1$ of the harmonic measure of $E$ in $D_{\rho_2}\setminus E$; it is positive because a segment has positive capacity and the circle is a compact subset of the slit disc, and it depends only on $\bar r$. Harmonic measure of $E$ increases with the outer domain, so the harmonic measure in $D_\rho\setminus E$ is at least $m_*$ on $|z|=\rho_1$; it vanishes on $|z|=\rho$. The maximum principle on the annulus $\rho_1<|z|<\rho$ then gives $\varpi\ge m_*\log(\rho/x)/\log(\rho/\rho_1)\ge m_*s/(4\log(1/\rho_1))$, since $\log(\rho/x)\ge\log(1+s/2)\ge s/4$. Hence $\log(1/\eta)\le4\log6\,\log(1/\rho_1)/(m_*s)$, and $1/s=\max\{4R^2,4/(1-\bar r)\}$.
\end{proof}

\begin{proof}[Proof of Theorem~\ref{thm:philog}]
\emph{Upper bound.} Take $c=(1,\dots,1)$ of length $T=\lceil\log(1/\delta)/\log(1/\bar r)\rceil\le L$. Then $g-p_c=r^T/(1-r)$, so $\psi(g-p_c)\le\bar r^T/(1-\bar r)\le\delta/(1-\bar r)$ and $\omega_L(\delta)\le\delta\sqrt T+B\delta/(1-\bar r)\le C\delta\sqrt{\log(1/\delta)}$ for small $\delta$.

\emph{Lower bound.} By part (iv) it suffices to bound $\delta\|c\|_2+B\psi(g-p_c)$ below for every $c$. Let $R=\|c\|_2$ and $\eta=\|g-p_c\|_\infty\le\psi(g-p_c)$. If $\eta\ge\delta\sqrt{\log(1/\delta)}$, the objective is at least $B\delta\sqrt{\log(1/\delta)}$. Otherwise $\eta<\sqrt\delta$, so $\log(1/\eta)>\frac12\log(1/\delta)$, and Lemma~\ref{lem:hardy} gives $K\max\{R^2,(1-\bar r)^{-1}\}>\frac12\log(1/\delta)$. Choose $\delta_0$ with $\log(1/\delta_0)\ge2K/(1-\bar r)$; then for $\delta\le\delta_0$ the maximum is $R^2$, so $R\ge\sqrt{\log(1/\delta)/(2K)}$ and the objective is at least $\delta\sqrt{\log(1/\delta)/(2K)}$. Hence $c=\min\{B,(2K)^{-1/2}\}$ works for every $L$.

\emph{Cost.} By part (iv), $N$ certifies $\Phi$ to half-width $\epsilon$ only if $\omega_L(2z_{1-\alpha}\sigma/\sqrt N)\le2\epsilon$, which with the lower bound forces $N\ge c'\sigma^2\epsilon^{-2}\log(1/\epsilon)$ for small $\epsilon$; the upper bound and part (iv) give a matching $N\le C'\sigma^2\epsilon^{-2}\log(1/\epsilon)$. On $\mathcal M(\bar r,B)$, $\Phi\le B/(1-\bar r)$, and when $a+\kappa B/(1-\bar r)<1$ the map $\Phi\mapsto S$ is increasing with derivative $\kappa^2/(1-a-\kappa\Phi)^2$ bounded above and below, so intervals for $\Phi$ and for $S$ convert into each other with constant changes in accuracy.
\end{proof}

\begin{corollary}[The rates of this thesis as cases of one picture]\label{cor:cases}
\begin{enumerate}[label=(\alph*)]
\item \emph{Safety with a memory bound.} $N_{\rm cert}(\epsilon;\Phi)\asymp\sigma^2\epsilon^{-2}\log(1/\epsilon)$, and the same for $S$ (Theorem~\ref{thm:philog}). Part (i) alone gives the upper bound $O(\sigma^2\epsilon^{-2}\log(1/\epsilon)/\log(1/\bar r))$, the rate of Theorem~\ref{p2:thm:certcost}; Theorem~\ref{thm:philog} shows that its logarithm is necessary. \hfill\Proved
\item \emph{Safety without a memory bound.} $g(r)=(1-r)^{-1}$ satisfies the condition of (iii): no finite number of experiments certifies the risk, as Theorem~\ref{p2:thm:window}(iii) shows for finite windows. \hfill\ProvedIn{CH.~3}
\item \emph{A single response tap, and discounted feedback.} For $g(r)=r^m$, $\omega_L(\delta)\le\delta$; for $g(r)=(1-\beta r)^{-1}$ with $\beta<1$, $\omega_\infty(\delta)\le\delta/\sqrt{1-\beta^2}$. Both are parametric (Proposition~\ref{prop:parametric}). \hfill\Proved
\item \emph{Identifying $K$ cohorts.} Recovering the atoms is not a linear question. It requires moments up to order $2K-1$ to be resolved against moment-matched alternatives whose higher moments differ by $O((\Delta\eta)^{2K-1})$, which costs $(\Delta\eta)^{-(4K-2)}$ experiments (Theorems~\ref{p1:thm:kcohort} and~\ref{p1:prop:esprit}). \hfill\ProvedIn{CH.~2}
\item \emph{Identifying a general memory measure.} Recovering $M$ in $W_1$ means certifying $F_g$ uniformly over all $1$-Lipschitz $g$. Their best approximations by degree-$T$ polynomials improve only as $1/T$, and the coefficients of those polynomials grow geometrically in $T$, so accuracy improves only logarithmically in $N$: the rate $1/\log N$ of Theorem~\ref{p1:thm:general}. \hfill\ProvedIn{CH.~2}
\end{enumerate}
\end{corollary}

The link in (e) between coefficient growth and the logarithmic rate is a heuristic reading of Chapter~\ref{ch:p1}'s proof through the dual form of part (iv), not a separate result here. The ladder is now exact at its first rungs: a single tap and any discounted feedback are parametric, the undiscounted total feedback that safety depends on costs exactly one extra logarithm, questions singular faster than $(1-r)^{-1/2}$ cannot be certified without a memory bound, and identification costs polynomially or logarithmically more.

\begin{remark}[Why coefficient energy is the wrong lower-bound quantity]\label{rem:energy}
Earlier drafts stated, as a target, a lower bound $N_{\rm cert}\ge c\,\sigma^2A_g(c'T_g(\epsilon))\epsilon^{-2}$ with constants uniform over a regular coefficient class. It is false. Fix $\epsilon_0>0$, $\beta\in(0,1)$, $\eta>0$ with $\eta/(1-\beta\bar r)<1/2$, and set $g_m(r)=(\epsilon_0/B)[(1-r/\bar r)^m+\eta/(1-\beta r)]$. Its coefficients are eventually positive with polynomial growth; on $[0,\bar r]$ it takes values in an interval of length below $3\epsilon_0/(2B)$, so a constant polynomial gives $\psi(g_m-p_c)<3\epsilon_0/(2B)$ with $\|c\|_2$ bounded and the question is certified to accuracy $\epsilon_0$ with no data. Yet $T_{g_m}(\epsilon_0)=m+1$ and $A_{g_m}(c'T_{g_m}(\epsilon_0))\epsilon_0^{-2}\ge B^{-2}(m/\bar r-\eta\beta)^2\to\infty$, because the alternating binomial tail $\sum_{t=T}^m(-1)^t\binom mt=(-1)^T\binom{m-1}{T-1}$ keeps $\tau_{g_m}(T)>\epsilon_0/2$ for $T\le m$. Taylor coefficients can be large while a different polynomial approximates the question cheaply; part (iv) measures the cheapest one.
\end{remark}

\paragraph{Numerical certificates.} The script \texttt{checks/check\_modulus.py} solves the primal (pairs of measures on a grid) and the dual (coefficient vectors) of part (iv) for the total feedback, for discounted feedback with $\beta=0.8$ and for the tap $r^3$, at $\bar r\in\{0.5,0.9\}$ and $\delta$ from $10^{-1}$ to $10^{-5}$, and stores both. Without any solver, the check rescales the stored measures to exact feasibility (a lower bound on $\omega_L$) and evaluates the stored coefficients over the whole interval with a curvature allowance between grid points (an upper bound, by weak duality). The brackets agree to a relative $4\times10^{-4}$ or better. Over four decades of $\delta$, $\omega_L(\delta)/\delta$ settles at the values $\|a\|_2$ of Proposition~\ref{prop:parametric} for the two parametric questions ($1/\sqrt{1-\beta^2}=5/3$ for the discounted question and $1$ for the tap) and keeps growing for $\Phi$ (from $6.6$ to $10.1$ at $\bar r=0.9$), while $\omega_L(\delta)/(\delta\sqrt{\log(1/\delta)})$ stays between $1.1$ and $4.4$. The check also verifies $\omega_L(0)=2B\,E_L(g)$ at small depths. These computations illustrate the theorems; they are not part of their proofs.

\section{Disclosure changes the contract, not the question}\label{sec:fw_disclosure}

Disclosure regimes are observation contracts. Class-tagged flow reveals each class's own moments, which turns identification into an ordinary estimation problem and is worth a factor $(\Delta\eta)^{-(4K-4)}$ in samples (Theorem~\ref{p1:thm:discvalue}). For safety the value of disclosure is different: it raises safe capacity exactly when it rules out the extremal compatible economy, and it is worth most when memory is long, where Theorem~\ref{thm:dich}(iii) says public data certify nothing (Theorem~\ref{p2:thm:disclosure}). Supervision that wants safety should therefore ask for the disclosures that bound memory or reveal $\Phi$, not for identities.

\section{A robust safety bound: nonlinear and changing agents}\label{sec:bl_robust}

Papers~1 and~2 work with agents whose rules are linear filters with fixed memory and weight. Real agents cap positions, trade only beyond thresholds, adapt their gains, and keep learning while the operator watches. This section and the next show that the safety side survives all of this under one condition: a bound on the agents' total sensitivity.


Return to the economy of Definition~\ref{def:cohort} and let everything about the agents vary over time.

\begin{assumption}[Bounded total sensitivity]\label{ass:sector}
For each cohort $i$ and period $t$, the memory rate $r_i(t)\in[0,1)$ is arbitrary, the weight satisfies $|q_i(t)|\le\bar q_i$, and the response rule satisfies the sector condition $|\varphi_{i,t}(y)|\le s_i|y|$ for all $y$. Write $\bar\Phi=\sum_i\bar q_i s_i$ for the total sensitivity.
\end{assumption}

The sector condition covers saturating momentum ($s_i=1$ for $\tanh$), position limits, threshold and dead-zone rules ($y\mathbf 1\{|y|>\tau\}$), and any rule whose response never exceeds a multiple of its signal. It does not require the rule to be smooth, monotone, stationary or known.

\begin{theorem}[Robust safety bound]\label{thm:sector}
Under Assumption~\ref{ass:sector}, if $|a|+\kappa\bar\Phi<1$ and the economy starts at rest, then for every shock sequence with $|w_t|\le\bar w$,
\[
\sup_{t\ge0}|x_t|\;\le\;\bar w\,\frac{\kappa}{1-|a|-\kappa\bar\Phi}.
\]
When $a\ge0$, this is the worst-case loss bound of the linear stationary economy with weights $\bar q_is_i$, so the summed-gain guard of Chapter~\ref{ch:p2} computed at total sensitivity $\bar\Phi$ is safe for every economy satisfying the assumption. \hfill\Proved
\end{theorem}

\begin{proof}
Let $Z_t=\max_{s\le t}|x_s|$. Each memory state is a convex combination of its previous value and the current state, so $|y_{i,t}|\le r_i(t)|y_{i,t-1}|+(1-r_i(t))|x_t|\le\max(|y_{i,t-1}|,|x_t|)$, and by induction from $y_{i,0}=0$, $|y_{i,t}|\le Z_t$. Then
\[
|x_{t+1}|\le|a|\,|x_t|+\kappa\Bigl(\sum_i|q_i(t)|\,|\varphi_{i,t}(y_{i,t})|+|w_{t+1}|\Bigr)\le(|a|+\kappa\bar\Phi)\,Z_t+\kappa\bar w .
\]
Let $Z^*=\kappa\bar w/(1-|a|-\kappa\bar\Phi)$. If $Z_t\le Z^*$ then $|x_{t+1}|\le(|a|+\kappa\bar\Phi)Z^*+\kappa\bar w=Z^*$, so $Z_{t+1}\le Z^*$; and $Z_0=0$. For the second claim, with $a\ge0$ the linear stationary economy with weights $\bar q_is_i$ is a positive system, its impulse response is non-negative, and its worst-case peak response to shocks bounded by $\bar w$ is $\bar w\sum_th_t=\bar w\,\kappa/(1-a-\kappa\bar\Phi)$ (equation~\eqref{eq:fw_S}). The bound of the theorem therefore coincides with the quantity the guard controls (Lemma~\ref{p2:thm:guard}).
\end{proof}

\paragraph{What the theorem says.} An operator who wants a safe admission level does not need to know how many agents there are, what their memories are, whether their rules are linear, or whether they change over time. It needs one number, an upper bound on the agents' total sensitivity $\bar\Phi$. The bound is tight: when the agents are linear, stationary, positive and at their maximal weights, it is attained. Nonlinearity, memory changes and negative feedback can only make the market safer than the bound says.

\begin{remark}[What the theorem does not cover]
Rules that amplify their signal, $|\varphi(y)|>s|y|$ for the declared $s$, break the bound; so do changes in the market's own persistence $a$ or impact $\kappa$, which the operator must certify separately (Chapter~\ref{ch:p2}, Section~\ref{p2:sec:infosafety}). The bound controls the peak of the price path, which is what a fixed position's loss depends on; dynamic positions are handled by the balance-sheet guard of Chapter~\ref{ch:p2}.
\end{remark}

\begin{proposition}[Several assets]\label{prop:multiasset}
Let $x_t\in\R^d$ evolve as $x_{t+1}=Ax_t+\mathcal K\bigl(\sum_iQ_i(t)\varphi_{i,t}(y_{i,t})+w_{t+1}\bigr)$, with memory states $y_{i,t}=R_i(t)y_{i,t-1}+(I-R_i(t))x_t$, $R_i(t)$ diagonal with entries in $[0,1)$, sector conditions $|\varphi_{i,t}(y)|\le s_i|y|$ componentwise, and $\|Q_i(t)\|_\infty\le\bar q_i$. If $\|A\|_\infty+\|\mathcal K\|_\infty\sum_i\bar q_is_i<1$, then $\sup_t\|x_t\|_\infty\le\bar w\,\|\mathcal K\|_\infty/(1-\|A\|_\infty-\|\mathcal K\|_\infty\sum_i\bar q_is_i)$ for $\|w_t\|_\infty\le\bar w$. \hfill\Proved
\end{proposition}

\begin{proof}
The proof of Theorem~\ref{thm:sector} goes through with absolute values replaced by the sup norm and scalars by induced sup norms: each component of $y_{i,t}$ is a convex combination of its past value and the corresponding component of $x_t$.
\end{proof}

The induced-norm bound discards the pattern of cross-impact. For positive cross-impact, the following comparison gives the sharp replacement.

\begin{proposition}[Positive cross-impact comparison]\label{prop:multiasset_positive}
Assume the economy starts at rest, $A,\mathcal K,\overline Q_i$ are entrywise non-negative, $|Q_i(t)|\le\overline Q_i$ entrywise, and $|\varphi_{i,t}(y)|\le S_i|y|$ for a non-negative diagonal matrix $S_i$. Memory matrices $R_i(t)$ are diagonal with entries in $[0,1)$. Put $G=A+\mathcal K\sum_i\overline Q_i S_i$. If $\rho(G)<1$, then for every $\|w_t\|_\infty\le\bar w$,
\[
\sup_t\|x_t\|_\infty\le
\bar w\,\bigl\|(I-G)^{-1}\mathcal K\bigr\|_\infty.
\]
This bound is sharp over the stated class. \hfill\Proved
\end{proposition}

\begin{proof}
The Neumann series $\sum_{j\ge0}G^j=(I-G)^{-1}$ has non-negative entries. Write $\mathbf1$ for the vector of ones and set $z=\bar w(I-G)^{-1}\mathcal K\mathbf1\ge0$. Then $Gz+\bar w\mathcal K\mathbf1=z$. Suppose $|x_t|\le z$ and $|y_{i,t}|\le z$ componentwise. The plant and the sector bounds give
\[
|x_{t+1}|\le A z+\mathcal K\Bigl(\sum_i\overline Q_i S_i z+\bar w\mathbf1\Bigr)=z.
\]
Each next memory coordinate is a convex combination of its previous value and the corresponding next state coordinate, so its absolute value is also bounded by the corresponding coordinate of $z$. Induction from rest proves the componentwise bound for all times. Because the resolvent times $\mathcal K$ is non-negative, its induced sup norm is its maximum row sum, and $\|z\|_\infty$ equals the displayed bound.

For sharpness, choose $R_i(t)=0$, $Q_i(t)=\overline Q_i$, $\varphi_{i,t}(y)=S_i y$, and $w_t=\bar w\mathbf1$. The state then satisfies $x_{t+1}=Gx_t+\bar w\mathcal K\mathbf1$, so from rest $x_t=\bar w\sum_{j<t}G^j\mathcal K\mathbf1$ increases componentwise to $z$. Its supremum norm therefore approaches the bound. Sharpness means a supremum over the admitted model class, not attainment at a finite time in every economy.
\end{proof}

\paragraph{Registered numerical validation.}
% Generated by checks/export_multiasset_numbers.py from R_multiasset_recovery.json
\newcommand{\MultiassetCases}{128}
\newcommand{\MultiassetPassed}{128}

The specification and its hash were pushed before the seeded instances were generated, and the unlock counter was advanced before the single opening. The joint primary required both componentwise enclosure for nonlinear, time-varying cohorts and convergence to the bound in the maximal positive linear instance. It passed in \MultiassetPassed{} of \MultiassetCases{} registered random instances. The generated matrices, tolerances and diagnostics are recorded in \texttt{checks/R\_multiasset.json}. This is a numerical check of the proved statement under the registered construction, not evidence for it.

\section{Certificates expire: learning agents}\label{sec:bl_learning}

Theorem~\ref{thm:sector} allows weights to change, provided they stay below their bounds. When agents learn, the bounds themselves are the issue: a certificate issued today says nothing about weights next month. The theorem turns this into a computable lifetime.

\begin{corollary}[Certificate lifetime]\label{thm:lifetime}
Suppose a certificate at time $0$ establishes $|q_i(0)|\le q_i^+$ for each cohort, and the guard's loss budget is met for every total sensitivity up to $\Phi_{\max}$, that is $\kappa\bar w/(1-a-\kappa\Phi_{\max})$ times the position equals the budget. If each weight changes by at most $\gamma_i$ per period, $|q_i(t+1)-q_i(t)|\le\gamma_i$, and $\sum_iq_i^+s_i\le\Phi_{\max}$ with $\sum_i\gamma_is_i>0$, then the guard remains safe for at least
\[
T_{\rm valid}=\Bigl\lfloor\frac{\Phi_{\max}-\sum_iq_i^+s_i}{\sum_i\gamma_is_i}\Bigr\rfloor
\]
periods under the sector bounds. \hfill\Proved
\end{corollary}

\begin{proof}
On $[0,T_{\rm valid}]$, $|q_i(t)|\le q_i^++\gamma_iT_{\rm valid}=:\bar q_i$ and $\sum_i\bar q_is_i\le\Phi_{\max}$. Theorem~\ref{thm:sector} applies on that horizon, since its induction uses only periods up to the current one.
\end{proof}

The lifetime is the certificate's safety margin divided by the agents' speed of learning. This is a guaranteed lower bound, not a claim that loss first crosses the budget at this time. Realized safety can last much longer if learning slows, changes sign, or stops below the sensitivity budget. When $\sum_i\gamma_i s_i=0$, the sensitivity envelope does not grow and the same certificate remains valid for every horizon under the stated sector and plant assumptions; the displayed quotient is not used. If the initial sensitivity bound already exceeds $\Phi_{\max}$, this argument supplies no valid initial certificate. Recertification frequency can therefore be based on the certified sensitivity margin and the learning-speed bound, while additional information about the realized learning path may justify a longer horizon.

\paragraph{Formal scope.}
The finite-cohort, finite-horizon safety induction is checked by \nolinkurl{CurrentMM.thesis_sector_safety}. The drift, floor and sensitivity calculations are checked by \nolinkurl{CurrentMM.thesis_lifetime_envelope} and its helper lemmas. The deterministic certification arithmetic is checked by \nolinkurl{CurrentMM.thesis_variance_radius} and \nolinkurl{CurrentMM.thesis_interval_half_width}; the new bias lemma covers finite positive atomic measures. For several assets, \nolinkurl{CurrentMM.thesis_multiasset_comparison} checks the comparison step of Proposition~\ref{prop:multiasset_positive} for any number of assets and cohorts, nonlinear time-varying sector-bounded responses, entrywise weight bounds and all times: every non-negative supersolution bounds every state and memory coordinate; \nolinkurl{CurrentMM.thesis_multiasset_linear_monotone} checks that the maximal positive linear economy rises from rest and stays below every supersolution. The Neumann series that produces the resolvent supersolution and the limit in the sharpness argument are written proofs. For finite atomic measures, \nolinkurl{CurrentMM.thesis_modulus_weak_duality} and \nolinkurl{CurrentMM.thesis_modulus_affine_bias} check the deterministic half of Theorem~\ref{thm:dich}(iv): the weak-duality bound on the modulus and the bias of the affine interval. These declarations do not formalize Gaussian coverage, a general-measure integration theorem, the minimax step and the lower bound of part (iv), or Theorem~\ref{thm:philog}, which rest on Sion's minimax theorem and on potential theory. The coverage register records these boundaries separately.

\paragraph{Identification under learning, by contrast.} Chapter~\ref{ch:p1}, Section~\ref{p1:sec:learners}, shows that trained online learners obey the same $\delta^{-4}$ law for separating memory rates when the competing populations' weights match, and that populations whose learned weights differ are told apart by those weights instead. Learning therefore leaves identification as hard as before while making certification only perishable, not harder per renewal.

\section{Markets of trading algorithms that learn}\label{sec:bl_ai}

Markets in which learning algorithms, including reinforcement-learning and language-model agents, trade with each other are where the results of this chapter matter most. Two of them apply directly. Certification of safety needs only a bound on total sensitivity (Theorem~\ref{thm:sector}); and the certificate lasts for its safety margin divided by the agents' learning speed (Corollary~\ref{thm:lifetime}). The second result gives a supervisory rule with no counterpart in identity-based regimes: recertify at a frequency set by how fast the agents update.

\paragraph{A registered simulation with learning agents.}
% generated by check_multiagent_v2.py; do not edit
\newcommand{\MAvTwoRuns}{200}
\newcommand{\MAvTwoPonePasses}{200}
\newcommand{\MAvTwoAdvRuns}{100}
\newcommand{\MAvTwoPtwoHits}{98}
\newcommand{\MAvTwoPtwoRate}{98.0\%}
\newcommand{\MAvTwoPtwoVerdict}{passed}

% generated by export_multiagent_v2_table.py; do not edit
\newcommand{\MAvTwoTableRows}{%
Adversarial & $0.001$ & $300$ & 25/25 & 1.27 & 25/25 & 3.97 & 2.7\\
Adversarial & $0.003$ & $100$ & 25/25 & 1.30 & 25/25 & 3.99 & 2.7\\
Adversarial & $0.010$ & $30$ & 25/25 & 1.33 & 25/25 & 4.03 & 2.9\\
Adversarial & $0.030$ & $10$ & 25/25 & 1.80 & 17/25 & 4.70 & 3.7\\
Random & $0.001$ & $300$ & 0/25 & --- & 0/25 & --- & 5.4\\
Random & $0.003$ & $100$ & 0/25 & --- & 0/25 & --- & 5.8\\
Random & $0.010$ & $30$ & 0/25 & --- & 0/25 & --- & 6.2\\
Random & $0.030$ & $10$ & 0/25 & --- & 0/25 & --- & 7.1\\}

Corollary~\ref{thm:lifetime} is a guarantee, so a simulation cannot confirm it; what a simulation can show is whether it is informative when agents learn from their own trading. We registered one (specification hash \texttt{58c32f9b\ldots}, pushed with its runner in a commit of its own before the run). The economy is that of Definition~\ref{def:cohort} with $a=0.3$, $\kappa=0.4$ and four cohorts with memory rates $0.2,0.5,0.7,0.85$, each following the saturating momentum rule $\varphi(y)=\tanh y$ (sector bound $1$). Each cohort learns its weight online by a gradient step on its own trading profit, clipped to the declared learning rate $\gamma$, so weights rise or fall but never move by more than $\gamma$ per period. Weights start at $0.05$ (total sensitivity $0.2$), the budget is $\Phi_{\max}=1.4$ (so the loss budget is $\kappa\bar w/(1-a-\kappa\Phi_{\max})=2.86$ for $\bar w=1$), and $T_{\rm valid}=\lfloor1.2/(4\gamma)\rfloor$ is $300$, $100$, $30$ and $10$ periods for $\gamma=0.001,0.003,0.01,0.03$. Shocks are bounded by $\bar w$: in the adversarial regime they push in the direction of the price displacement nine times in ten; in the random regime they are uniform. There are 25 runs for each regime and learning rate, each lasting five lifetimes.

The registered check (P1) held in \MAvTwoPonePasses{} of \MAvTwoRuns{} runs: up to $T_{\rm valid}$ every price path stayed within the loss budget and total sensitivity within $\Phi_{\max}$, as the corollary guarantees. The informative primary (P2) asked whether learning actually exhausts the sensitivity budget within three lifetimes in at least half of the adversarial runs. It \MAvTwoPtwoVerdict{}: it did in \MAvTwoPtwoHits{} of \MAvTwoAdvRuns{} runs. Table~\ref{tab:multiagent_v2} shows that under adversarial shocks the budget is exhausted after 1.3 to 1.8 lifetimes, and the realized loss first exceeds its budget after four to five. So the lifetime bound is close to tight when trend-following pays, and a certificate that is not renewed does eventually admit losses beyond its budget. Under random shocks trend-following loses money on average, weights decay towards zero and no run comes near the budget: there the bound is conservative, because it charges every period for the fastest possible learning.

\begin{table}[ht]
\centering\small
\begin{threeparttable}
\caption{Registered simulation with learning agents: certificate lifetime and realized loss}\label{tab:multiagent_v2}
\begin{tabular}{lccccccc}
\toprule
Shocks & $\gamma$ & $T_{\rm valid}$ & Budget exhausted & Median time & Loss breached & Median time & Median margin\\
\midrule
\MAvTwoTableRows
\bottomrule
\end{tabular}
\begin{tablenotes}\footnotesize\item 25 runs per row, horizon $5T_{\rm valid}$. ``Budget exhausted'': runs in which total sensitivity exceeds $\Phi_{\max}$; ``Loss breached'': runs in which $|x_t|$ exceeds the loss budget. Median times are in units of $T_{\rm valid}$, over the runs in which the event occurs. ``Median margin'': loss budget divided by the largest $|x_t|$ up to $T_{\rm valid}$. Source: \texttt{checks/R\_multiagent\_v2.json}. An earlier simulation, whose weight updates could only increase and so could not test anything, is withdrawn (Chapter~\ref{ch:conclusion}).\end{tablenotes}
\end{threeparttable}
\end{table}

\section{Agents who respond to the guard}\label{sec:bl_strategic}

A guard changes the environment it certifies. If agents learn that the operator throttles admission when measured sensitivity is high, they may change how they trade.

\begin{conjecture}[Guard-consistent equilibrium]\label{conj:strategic}
Suppose each cohort chooses its weight $q_i\in[0,\bar q_i]$ to maximize expected trading profit net of a convex cost, given an admission rule that throttles the dealer's position when certified total sensitivity exceeds $\Phi_{\max}$, and given the equilibrium reversion supplied by dealers (Proposition~\ref{p2:prop:equil}). Then (i) an equilibrium exists; (ii) in every equilibrium the guard is safe; and (iii) equilibrium total sensitivity is weakly lower than without the guard. \hfill\Conj
\end{conjecture}

\begin{remark}[Counterexample to existence under a discontinuous guard]\label{rem:strategic_counterexample}
Part (i) of the conjecture is false under the stated assumptions. Consider one cohort with $q\in[0,1]$, a perfect certificate $\Phi=q$, and admitted dealer position $I(q)=1$ for $q\le1/2$ and $I(q)=0$ for $q>1/2$. This is an admission rule that throttles above the certified threshold. Let expected trading profit, as a function of weight and admitted position, be
\[
\pi(q,I)=(1-I)(1-q),\qquad c(q)=0.
\]
The profit primitive is continuous in $(q,I)$ and the zero cost is convex. After composition with the guard, utility is zero at $q\le1/2$ and equals $1-q$ at $q>1/2$. If $q\le1/2$, the admissible deviation $q'=3/4$ strictly increases utility. If $q>1/2$, the admissible deviation $q'=(q+1/2)/2$ lies strictly between $1/2$ and $q$ and also strictly increases utility. Thus no weight maximizes utility: the one-player game has no pure equilibrium despite its compact strategy set, convex cost, and exact sensitivity certificate.

Randomization does not restore existence. Utility is strictly below $1/2$ everywhere, so its expectation under any probability distribution is strictly below $1/2$; a pure weight sufficiently close to the threshold from above earns more. The universal profitable-deviation statement and nonexistence of a pure maximizer are formalized as \nolinkurl{CurrentMM.thesis_strategic_profitable_deviation} and \nolinkurl{CurrentMM.thesis_strategic_no_equilibrium}. The probability-integral step for arbitrary mixed strategies is a written argument, not part of those Lean declarations. Exact rational diagnostics are recorded separately in \texttt{checks/R\_strategic.json}.

The failure concerns equilibrium existence, not the safety invariant. For example take $a=0$, $\kappa=1/4$, and $\bar w=1$. On admitted profiles $q\le1/2$, unit position loss is bounded by $2/7$; on throttled profiles the position is zero. Every admitted decision therefore respects this budget, although the agents' game has no equilibrium. A revised strategic theorem needs assumptions that rule out the best-response discontinuity. Convex cost alone also supplies no ordering of equilibrium sensitivity with and without a guard; that comparative-static claim remains open under suitable additional payoff restrictions.
\end{remark}

\section{What the framework does not cover}\label{sec:fw_scope}

The framework is linear in the observable layer and assumes the reset contract for its sharpest statements. Sections~\ref{sec:bl_robust} and~\ref{sec:bl_learning} extend the safety results to sector-bounded nonlinear and to learning agents; strategic equilibrium existence fails without additional assumptions (Remark~\ref{rem:strategic_counterexample}); existence for suitably regularized guards remains open. Passive observational data are slower but obey the same polynomial order for identification (Theorem~\ref{p1:prop:passive}); a finite-sample theory of certification from observational order-book data, with dependence and non-stationarity, is not developed here, and the empirical certificates of Chapter~\ref{ch:p2} are bootstrap diagnostics rather than guarantees.
